\documentclass[12pt]{article}
\usepackage[T1]{fontenc}
\usepackage[utf8]{inputenc}
\usepackage{lmodern}
\usepackage{textcomp}
\usepackage{enumitem}
\usepackage{geometry}
\geometry{margin=1in}
\begin{document}
\begin{center}
Answer Key - Mathematics - K-12 Level Assessment 10 \
\small (Answers are ordered exactly as in the question paper.)
\end{center}

\section*{Multiple Choice}
\begin{enumerate}[label=\arabic*.]
\item \textbf{Answer:} B
\\ \textit{Options:} A) 6; B) 18; C) 36; D) 180
\\ \textit{Note:} The greatest common factor of 36 and 90 is 18 because it is the largest integer that divides both numbers without leaving a remainder, as evidenced by their prime factorizations: 36 is $2^2$ × $3^2$ and 90 is 2 × $3^2$ × 5, where the common factors are 2 and 3, leading to 2 × 3 × 3 = 18.
\item \textbf{Answer:} B
\\ \textit{Options:} A) The standard deviation of the population is also zero.; B) The sample mean and sample median are equal.; C) The sample may have outliers.; D) The population has a symmetric distribution.
\\ \textit{Note:} When the standard deviation is zero, it means all values in the sample are identical, leading to the sample mean and sample median being equal because both measures represent the same value.
\item \textbf{Answer:} C
\\ \textit{Options:} A) 34; B) 2; C) 17; D) 3
\\ \textit{Note:} The least possible positive integer value of \( n \) is 17 because it is necessary for \( n \) to contain the prime factor 17 to ensure that the product \( 18 \cdot n \cdot 34 \) is a perfect square, as both 18 and 34 contribute prime factors that require balancing to form a complete set of pairs.
\item \textbf{Answer:} A
\\ \textit{Options:} A) 29; B) 28; C) 27; D) 26
\\ \textit{Note:} The correct answer is 29 because when substituting the given values into the expression for M, we find that j = 51 and k = -11, leading to j + 2 k = 51 + 2(-11) = 29.
\item \textbf{Answer:} B
\\ \textit{Options:} A) -9; B) -14; C) 14; D) 9
\\ \textit{Note:} The sum of all integer solutions to the inequality |n| \textless{} |n-3| \textless{} 9 yields the integer solutions n = -8, -7, -6, -5, -4, -3, -2, -1, 0, 1, 2, 3, 4, 5, 6, 7, 8, which sum to -14.
\end{enumerate}

\section*{True / False}
\begin{enumerate}[label=\arabic*.]
\setcounter{enumi}{5}
\item \textbf{Answer:} False
\\ \textit{Options:} A) True; B) False
\\ \textit{Note:} The statement is false because if f(2) = 5, it does not imply that f(-2) will also equal 5, as the polynomial's behavior can differ for negative and positive inputs.
\item \textbf{Answer:} False
\\ \textit{Options:} A) True; B) False
\\ \textit{Note:} A prism has at least five faces, including two parallel bases, and cannot roll like a cylinder due to its angular edges and flat surfaces.
\item \textbf{Answer:} False
\\ \textit{Options:} A) True; B) False
\\ \textit{Note:} The statement is false because if \( x \) is twice \( y \), their sum cannot equal 42 while also being inversely proportional, as inversely proportional quantities cannot maintain a fixed ratio.
\item \textbf{Answer:} True
\\ \textit{Options:} A) True; B) False
\\ \textit{Note:} The statement is true because solving the equation $3^(x$ - 3) + 10 = 19 leads to $3^(x$ - 3) = 9, and since 9 can be expressed as $3^2$, it follows that x - 3 = 2, giving x = 5.
\item \textbf{Answer:} True
\\ \textit{Options:} A) True; B) False
\\ \textit{Note:} The statement is true because when you add 17 to both sides of the equation 29 = x - 17, you find that x equals 46.
\end{enumerate}

\section*{Long Form Response}
\begin{enumerate}[label=\arabic*.]
\setcounter{enumi}{10}
\item \textbf{Answer:} We seek to use Vieta's formulas. To be able to apply the formulas, we multiply both sides by $(x-1)(x-5)(x-10)(x-25)$ to eliminate the fractions. This gives \[\begin{aligned}&\quad (x-5)(x-10)(x-25) + (x-1)(x-10)(x-25) \\& + (x-1)(x-5)(x-25) + (x-1)(x-5)(x-10) = 2(x-1)(x-5)(x-10)(x-25). \end{aligned}\](Be careful! We may have introduced one of the roots $x = 1, 5, 10, 25$ into this equation when we multiplied by $(x-1)(x-5)(x-10)(x-25).$ However, note that none of $x = 1, 5, 10, 25$ satisfy our new equation, since plugging each one in gives the false equation $1=0.$ Therefore, the roots of this new polynomial equation are the same as the roots of the original equation, and we may proceed.) The left-hand side has degree $3$ while the right-hand side has degree $4,$ so when we move all the terms to the right-hand side, we will have a $4$th degree polynomial equation. To find the sum of the roots, it suffices to know the coefficients of $x^4$ and $x^3.$ The coefficient of $x^4$ on the right-hand side is $2,$ while the coefficients of $x^3$ on the left-hand and right-hand sides are $4$ and $2(-1-5-10-25) = -82,$ respectively. Therefore, when we move all the terms to the right-hand side, the resulting equation will be of the form \[ 0 = 2 x^4 - 86 x^3 + \cdots\]It follows that the sum of the roots is $\tfrac{86}{2} = \boxed{43}.$
\\ \textit{Note:} correct because it applies Vieta's formulas after eliminating the fractions by multiplying through by the common denominator, while also acknowledging that the roots introduced by this multiplication do not satisfy the original equation, ensuring that the solution remains valid.
\item \textbf{Answer:} To find the percentage of Yum-Yum Good candy bars containing more than 225 calories, we first need to convert 225 calories into a z-score using the formula: z = (X - μ) / σ, where X is 225, μ is the mean (210), and σ is the standard deviation (10). This gives us a z-score of (225 - 210) / 10 = 1.5. Using standard normal distribution tables or a calculator, we find that the area to the left of a z-score of 1.5 is approximately 0.9332. Therefore, the percentage of candy bars containing more than 225 calories is 1 - 0.9332, which equals about 0.0668 or 6.68\%. This calculation shows that only a small portion of the candy bars exceed 225 calories.
\\ \textit{Note:} correct because it accurately calculates the z-score corresponding to 225 calories and uses the standard normal distribution to determine the percentage of candy bars that exceed this calorie amount by finding the area to the left of the z-score and subtracting it from 1.
\end{enumerate}

\vfill
\noindent\textit{End of Answer Key}
\end{document}